\documentclass[border=4pt]{standalone}
\usepackage{amsmath,amssymb,mathtools,xcolor,varwidth}
\begin{document}
\begin{varwidth}{28em}
\large\bfseries
In the same figure as before (Pl.~1, Fig.~6), the curve $acE$ stands above the base $AE$, but now we take \emph{unequal} bases $AB$, $BC$, $CD$, $\&c.$, the first wide and those following narrower. Upon them we erect the stepped parallelograms under the curve, just as in \textbf{Lemma~II}, except that here their breadths need not agree. The earlier lemma proved the ultimate equality of inscribed and circumscribed figures only for equal bases; this lemma asserts that the same conclusion survives when the breadths differ. We must therefore show that the total excess of the circumscribed figure over the inscribed one still vanishes as every base shrinks. The whole of that excess will be trapped within a single rectangle built on the \emph{widest} base.
\end{varwidth}
\end{document}
